\section{Introduction}
\label{sec:intro}

Retrieval-Augmented Generation (RAG) systems have been developed to enhance large language models (LLMs) by integrating external knowledge sources~\cite{sudhi2024rag, es2024ragas, salemi2024evaluating}. This innovative integration allows LLMs to generate more accurate and contextually relevant responses, significantly improving their utility in real-world applications. By adapting to specific domain knowledge~\cite{tu2024r}, RAG systems ensure that the information provided is not only pertinent but also tailored to the user's needs. Furthermore, they offer access to up-to-date information~\cite{zhao2024retrieval}, which is crucial in rapidly evolving fields. Chunking plays a vital role in facilitating the retrieval-augmented generation process~\cite{lyu2024crud}. By breaking down a large external text corpus into smaller, more manageable segments, chunking significantly enhances the accuracy of information retrieval. This approach allows for more targeted similarity searches, ensuring that the retrieved content is directly relevant to user queries.

However, existing RAG systems have key limitations that hinder their performance. \textbf{First}, many methods rely on flat data representations, restricting their ability to understand and retrieve information based on intricate relationships between entities. \textbf{Second}, these systems often lack the contextual awareness needed to maintain coherence across various entities and their interrelations, resulting in responses that may not fully address user queries. For example, consider a user asking, ``How does the rise of electric vehicles influence urban air quality and public transportation infrastructure?'' Existing RAG methods might retrieve separate documents on electric vehicles, air pollution, and public transportation challenges but struggle to synthesize this information into a cohesive response. They may fail to explain how the adoption of electric vehicles can improve air quality, which in turn could affect public transportation planning. As a result, the user may receive a fragmented answer that does not adequately capture the complex inter-dependencies among these topics.

To address these limitations, we propose incorporating graph structures into text indexing and relevant information retrieval. Graphs are particularly effective at representing the interdependencies among different entities~\cite{rampavsek2022recipe}, which enables a more nuanced understanding of relationships. The integration of graph-based knowledge structures facilitates the synthesis of information from multiple sources into coherent and contextually rich responses. Despite these advantages, developing a fast and scalable graph-empowered RAG system that efficiently handles varying query volumes is crucial. In this work, we achieve an effective and efficient RAG system by addressing three key challenges: i) \textbf{Comprehensive Information Retrieval}. Ensuring comprehensive information retrieval that captures the full context of inter-dependent entities from all documents; ii) \textbf{Enhanced Retrieval Efficiency}. Improving retrieval efficiency over the graph-based knowledge structures to significantly reduce response times; iii) \textbf{Rapid Adaptation to New Data}. Enabling quick adaptation to new data updates, ensuring the system remains relevant in dynamic environments.

In response to the outlined challenges, we propose \model, a model that seamlessly integrates a graph-based text indexing paradigm with a dual-level retrieval framework. This innovative approach enhances the system's capacity to capture complex inter-dependencies among entities, resulting in more coherent and contextually rich responses. \model\ employs efficient dual-level retrieval strategies: low-level retrieval, which focuses on precise information about specific entities and their relationships, and high-level retrieval, which encompasses broader topics and themes. By combining both detailed and conceptual retrieval, \model\ effectively accommodates a diverse range of quries, ensuring that users receive relevant and comprehensive responses tailored to their specific needs. Additionally, by integrating graph structures with vector representations, our framework facilitates efficient retrieval of related entities and relations while enhancing the comprehensiveness of results through relevant structural information from the constructed knowledge graph.

In summary, the key contributions of this work are highlighted as follows:

\begin{itemize}[leftmargin=*]

\item \textbf{General Aspect}. We emphasize the importance of developing a graph-empowered RAG system to overcome the limitations of existing methods. By integrating graph structures into text indexing, we can effectively represent complex interdependencies among entities, fostering a nuanced understanding of relationships and enabling coherent, contextually rich responses.

\item \textbf{Methodologies}. To enable an efficient and adaptive RAG system, we propose \model, which integrates a dual-level retrieval paradigm with graph-enhanced text indexing. This approach captures both low-level and high-level information for comprehensive, cost-effective retrieval. By eliminating the need to rebuild the entire index, \model\ reduces computational costs and accelerates adaptation, while its incremental update algorithm ensures timely integration of new data, maintaining effectiveness in dynamic environments.

\item \textbf{Experimental Findings}. Extensive experiments were conducted to evaluate the effectiveness of \model\ in comparison to existing RAG models. These assessments focused on several key dimensions, including retrieval accuracy, model ablation, response efficiency, and adaptability to new information. The results demonstrated significant improvements over baseline methods.

\end{itemize}